% 国赛论文 starter（cumcmthesis v2.7 clean，2026 年修订稿规范）
% 骨架冻结：不动 cumcmthesis.cls，不动边距/字体/标题样式。版式疼靠删内容解决。
%
% 编译（必须 XeLaTeX，至少两遍；推荐 latexmk）：
%   纸质版（含承诺书/编号专用页，打印装订用）：
%     latexmk -xelatex cumcm-paper
%   电子版（第一页必须为摘要页，单文件 PDF 提交用）：
%     将下一行换成 \documentclass[withoutpreface,bwprint]{cumcmthesis} 后再编
%
% 字体（开箱即用，无需手动装字体）：
%   cls 调用 Times New Roman / Arial + simkai.ttf / simsun.ttc（Windows 自带）。
%   Linux/macOS 无官方两文件时，下面 cumcm-fonts.sty 自动回退到系统现有中文字体
%   （Songti SC / Noto Serif CJK SC / Fandol），直接编译。要官方原字形才需跑一次
%   assets/scripts/setup-fonts.py（--copy-from <含两文件的目录> --dest userfonts）。
%   用稿时把 assets/templates/cumcmthesis/ 整个目录随稿走（cls + sty 不拆）。
%   字体替换警告必须认领。
\documentclass{cumcmthesis}
% 电子版提交时改用（去掉承诺书与编号专用页）：
% \documentclass[withoutpreface,bwprint]{cumcmthesis}
\usepackage{cumcm-fonts} % 官方字体缺失时自动回退；有官方字体时保持原样
%
% 标签命名空间：fig: / tab: / eq: / sec: ；交叉引用只用 \cref，不手写数字。
%   display 公式被引用才编号；图宽只用 \linewidth 分数，不用绝对英寸。
\documentclass{cumcmthesis}
% 电子版提交时改用（去掉承诺书与编号专用页）：
% \documentclass[withoutpreface,bwprint]{cumcmthesis}

\title{TODO：论文题目}
\tihao{TODO}                    % 题号：A/B/C/D/E 四选一填字母
\baominghao{TODO}               % 12 位报名参赛队号
\schoolname{TODO}               % 学校全称（电子版构建前必须清空匿名，见文末注释）
\membera{TODO}
\memberb{TODO}
\memberc{TODO}
\supervisor{TODO}
\yearinput{2026}
\monthinput{TODO}
\dayinput{TODO}

\begin{document}

\maketitle

\begin{abstract}
TODO：摘要。研究对象、主要方法、关键结果与结论，原则上不超过一页。
摘要页起以阿拉伯数字从 1 连续编页码（cls 自动处理，不手动改）。

\keywords{TODO\quad TODO\quad TODO}
\end{abstract}

% 2026 年规范：正文不要目录。下面一行保持注释。
% \tableofcontents

\section{问题重述}
\label{sec:restate}
TODO：准确、简洁重述赛题，不抄原题全文。

\section{问题分析}
\label{sec:analysis}
TODO：研究思路、问题分解、总体技术路线。

\section{模型假设}
\label{sec:assumptions}
\begin{enumerate}
    \item TODO：假设一；
    \item TODO：假设二。
\end{enumerate}

\section{符号说明}
\label{sec:symbols}
\begin{table}[htbp]
    \centering
    \caption{主要符号说明}
    \label{tab:symbols}
    \begin{tabular}{cc}
        \toprule[0.8pt]
        符号 & 含义 \\
        \midrule[0.5pt]
        $x$ & TODO：示例变量 \\
        \bottomrule[0.8pt]
    \end{tabular}
\end{table}

\section{模型建立与求解}
\label{sec:model}
TODO：模型、算法、求解过程和结果。引用示例如\cref{tab:symbols}、式\cref{eq:demo}。

\begin{equation}
E=mc^2
\label{eq:demo}
\end{equation}

% 插图示例（文件放 figures/，相对宽度）：
% \begin{figure}[htbp]
%     \centering
%     \includegraphics[width=.6\linewidth]{TODO-filename}
%     \caption{TODO：图题}
%     \label{fig:demo}
% \end{figure}

\section{模型检验与分析}
\label{sec:validation}
TODO：误差分析、敏感性分析、稳健性检验或其他验证。

\section{模型评价与推广}
\label{sec:evaluation}
TODO：优点、局限、可推广方向。

\begin{thebibliography}{9}
    \bibitem{ref1}
    TODO：作者.
    \newblock 文献题目.
    \newblock 出版信息, 年份.
\end{thebibliography}
% 文内引用资料（含网上资料）必须 \cite 标注 + 文末列出。无引用时删除此注释并确认。

\section*{人工智能工具使用声明}
TODO：依据竞赛期间有效的全国组委会规定填写；未使用时也应明确声明。

\begin{appendices}

\section{支撑材料文件列表}
TODO：逐项列出文件名、格式及用途。无支撑材料时写“本论文没有支撑材料”。

\section{源程序}
TODO：附全部完整、可运行源程序。无程序时写“本论文没有用到程序”。
% \begin{lstlisting}[language=matlab]
% TODO：源程序全文
% \end{lstlisting}

\end{appendices}

\end{document}

% ===== 电子版匿名构建前检查（提交前逐项划掉） =====
% - \schoolname / \membera / \memberb / \memberc / \supervisor 已清空或改为匿名
% - 摘要/正文/附录全文 grep 学校、姓名、赛区零命中
% - \documentclass 已换 [withoutpreface,bwprint]，第一页为摘要页
